\section{Conclusions}

The PSM voltage boundary is a real ellipse only after three facts have been
proved: the current-to-voltage map has full rank, its Gram matrix is positive
definite, and completing the square leaves the positive level
$U_{\max}^2$.  This route also exposes the exceptional zero-speed,
zero-resistance case.

The matrix $\Hpsm$ determines shape and orientation, the affine term determines
the center, and $U_{\max}$ determines scale.  Eigenvectors are axes because
they remove cross coupling; eigenvalues set inverse-squared axis lengths.
The classical discriminant and double-angle formula are coordinate expressions
of the same determinant and eigenvector facts.

Adding current and torque geometry turns operating strategies into contact
problems between level sets.  MTPA, field weakening, and MTPV are therefore not
isolated recipes but different active-constraint regimes of one geometric
optimization problem.  This viewpoint extends naturally to the extra
excitation dimension of an EESM, where rank rather than visual expectation
becomes decisive.
